%% Worksheets, self checks and unit test, AP Physics 1 Unit 2 (Force and Translational Dynamics), matching the guided notes 2.1 to 2.9.
%% Use g = 10 m/s^2 and ignore air resistance unless stated.
\renewcommand{\nocalc}{}
\newcommand{\gr}[2]{\begin{center}\begin{tikzpicture}\begin{axis}[gridaxis,#1]#2\end{axis}\end{tikzpicture}\end{center}}
\newcommand{\work}[1]{\ifkey\else\par\vspace{#1}\fi}
\newcommand{\ms}{\,\mathrm{m/s}}
\newcommand{\mss}{\,\mathrm{m/s^2}}
\newcommand{\m}{\,\mathrm{m}}
\newcommand{\s}{\,\mathrm{s}}
\newcommand{\N}{\,\mathrm{N}}
\newcommand{\kg}{\,\mathrm{kg}}
\newcommand{\Nm}{\,\mathrm{N/m}}

%%DOC 2.1A | Systems and Center of Mass
\begin{exo}{}
In physics, a system is
\choices{only the objects that are moving}{the object or objects we choose to study}{always a single object}{only an object that cannot bend}
\answer{B}{We choose the system; everything outside it is the environment.}
\end{exo}
\begin{exo}{}
A $3\kg$ mass is at $x=0$ and a $1\kg$ mass is at $x=4\m$. Where is their center of mass?
\choices{$x=2\m$}{$x=3\m$}{$x=1\m$}{$x=0.5\m$}
\answer{C}{$x_\text{cm}=\dfrac{3(0)+1(4)}{3+1}=1\m$, closer to the heavier mass.}
\end{exo}
\begin{exo}{}
Two equal masses are at $x=2\m$ and $x=8\m$. Where is their center of mass?
\choices{$x=5\m$}{$x=6\m$}{$x=4\m$}{$x=10\m$}
\answer{A}{Equal masses: the center of mass is halfway, at $5\m$.}
\end{exo}
\begin{exo}{}
Which forces can change the velocity of the center of mass of a system?
\choices{Internal forces only}{External forces only}{Both internal and external forces}{Neither}
\answer{B}{Internal forces come in third-law pairs inside the system and cancel for the system as a whole.}
\end{exo}
\begin{exo}{}
A firework moving along a parabola explodes into pieces in mid-air (no air resistance). After the explosion, the center of mass of the pieces
\choices{stops}{falls straight down}{moves straight up}{continues along the same parabola}
\answer{D}{The explosion is caused by internal forces; the only external force is still gravity.}
\end{exo}
\begin{exo}{}
A uniform $0.2\kg$ meter stick has a $0.2\kg$ mass attached at the $100\,\mathrm{cm}$ mark. Where is the center of mass of the stick and mass?
\choices{$50\,\mathrm{cm}$}{$75\,\mathrm{cm}$}{$100\,\mathrm{cm}$}{$62.5\,\mathrm{cm}$}
\answer{B}{The stick acts as $0.2\kg$ at $50\,\mathrm{cm}$: $\dfrac{0.2(50)+0.2(100)}{0.4}=75\,\mathrm{cm}$.}
\end{exo}
\begin{exo}{}
Two skaters at rest on frictionless ice push off each other. What happens to the center of mass of the two skaters?
\choices{It stays at rest.}{It moves toward the heavier skater.}{It moves toward the lighter skater.}{It accelerates.}
\answer{A}{The pushes are internal forces, and the net external force is zero.}
\end{exo}
\begin{exo}{}
When is it correct to model a car as a single object located at its center of mass?
\choices{When the car spins}{When its wheels deform}{When only the motion of the car as a whole matters}{Never}
\answer{C}{If the internal structure (rotation, deformation) does not affect the question, the object model works.}
\end{exo}
\Needspace{12\baselineskip}
\begin{exo}[Free response]{}
Three small masses are placed on an $xy$-grid: $1\kg$ at $(0,0)$, $2\kg$ at $(3\m,0)$ and $1\kg$ at $(0,4\m)$.
\begin{enumerate}[label=(\alph*)]
\item Find the coordinates of the center of mass.
\item The $2\kg$ mass is moved to $(3\m,4\m)$. Find the new center of mass.
\item The masses are held together by light rods. Explain why pushes and pulls between the masses cannot make the center of mass move.
\end{enumerate}
\work{4.5cm}
\begin{solution}
(a) $x=\frac{0+6+0}{4}=1.5\m$, $y=\frac{0+0+4}{4}=1\m$. \;(b) $x=1.5\m$, $y=\frac{0+8+4}{4}=3\m$. \;(c) These are internal forces: they come in equal and opposite pairs, so the net force on the system, and the acceleration of its center of mass, are unchanged.
\end{solution}
\end{exo}

%%DOC 2.2A | Forces and Free-Body Diagrams
\begin{exo}{}
What is the SI unit of force?
\choices{Kilogram}{Newton}{Joule}{Meter per second squared}
\answer{B}{$1\N=1\kg\cdot\mathrm{m/s^2}$.}
\end{exo}
\begin{exo}{}
What is the weight of a $5\kg$ mass near Earth's surface?
\choices{$5\N$}{$0.5\N$}{$50\N$}{$500\N$}
\answer{C}{$F_g=mg=5(10)=50\N$.}
\end{exo}
\begin{exo}{}
Which of these is a field (non-contact) force?
\choices{Friction}{Normal force}{Tension}{Gravitational force}
\answer{D}{Gravity acts at a distance; the others need contact.}
\end{exo}
\begin{exo}{}
A book rests on a table. Which forces act on the book?
\choices{Weight and normal force}{Weight only}{Normal force only}{Weight, normal force and friction}
\answer{A}{Earth pulls it down, the table pushes it up. With no sideways push there is no friction.}
\end{exo}
\begin{exo}{}
A free-body diagram should show
\choices{the forces the object exerts on other objects}{the velocity of the object}{only the forces exerted on the object}{the acceleration of the object}
\answer{C}{Only forces on the chosen object, as arrows starting at the dot.}
\end{exo}
\begin{exo}{}
A $3\kg$ lamp hangs at rest from a cord. What is the tension in the cord?
\choices{$3\N$}{$30\N$}{$15\N$}{$60\N$}
\answer{B}{At rest: $F_T=mg=30\N$.}
\end{exo}
\begin{exo}{}
A ball is thrown straight up. While it rises (no air resistance), which forces act on it?
\choices{Only gravity}{An upward force from the throw, and gravity}{Only an upward force}{No force}
\answer{A}{Once the ball leaves the hand, the hand no longer pushes it. Only gravity acts.}
\end{exo}
\begin{exo}{}
A rope pulls a box to the right across a rough floor at constant speed. Which force on the box points to the left?
\choices{The normal force}{The tension}{The weight}{Kinetic friction}
\answer{D}{Friction opposes the sliding motion.}
\end{exo}
\Needspace{14\baselineskip}
\begin{exo}[Free response]{}
Draw and label a free-body diagram for each situation. Make the relative lengths of the arrows reasonable.
\begin{enumerate}[label=(\alph*)]
\item A box slides to the right on a rough floor after being pushed (the push has stopped).
\item A skydiver falls at a constant speed (terminal velocity).
\item A box is pulled up a rough incline at constant speed by a rope parallel to the incline.
\end{enumerate}
\work{5.5cm}
\begin{solution}
(a) Weight down, normal force up (equal), kinetic friction to the left. No force to the right. \;(b) Weight down, air resistance up, equal in length (net force zero). \;(c) Weight straight down, normal force perpendicular to the incline, tension up the incline, kinetic friction down the incline; tension $=$ friction $+\;mg\sin\theta$.
\end{solution}
\end{exo}

%%DOC 2.3A | Newton's Third Law
\begin{exo}{}
A hammer hits a nail with a force of $500\N$. The force of the nail on the hammer is
\choices{zero}{less than $500\N$}{$500\N$, in the opposite direction}{more than $500\N$}
\answer{C}{Third law: equal magnitude, opposite direction.}
\end{exo}
\begin{exo}{}
A heavy truck collides with a small car. Compare the forces they exert on each other.
\choices{They have equal magnitudes.}{The truck exerts the larger force.}{The car exerts the larger force.}{The forces are zero.}
\answer{A}{Third-law pair: always equal in magnitude, whatever the masses.}
\end{exo}
\begin{exo}{}
In the same collision, which vehicle has the larger acceleration?
\choices{The truck}{The car}{They are equal.}{Neither accelerates.}
\answer{B}{Same force, smaller mass: $a=F/m$ is larger for the car.}
\end{exo}
\begin{exo}{}
Earth pulls a book downward. What is the third-law partner of this force?
\choices{The normal force of the table on the book}{The book pushing on the table}{The table pushing on the floor}{The book pulling Earth upward}
\answer{D}{The partner involves the same two objects (Earth and book) and the same type of force (gravity).}
\end{exo}
\begin{exo}{}
Why don't third-law force pairs cancel each other?
\choices{They have different magnitudes.}{They act on different objects.}{They act on the same object.}{They only exist for moving objects.}
\answer{B}{Forces cancel only when they act on the same object.}
\end{exo}
\begin{exo}{}
A swimmer pushes the water backward. The water pushes the swimmer
\choices{forward}{backward}{upward}{not at all}
\answer{A}{The reaction force of the water drives the swimmer forward.}
\end{exo}
\begin{exo}{}
Earth pulls a $60\kg$ student with a force of $600\N$. With what force does the student pull Earth?
\choices{$0$}{$6\N$}{$600\N$}{$6000\N$}
\answer{C}{Equal magnitude, opposite direction. Earth's acceleration is tiny because its mass is huge.}
\end{exo}
\begin{exo}{}
A horse pulls a cart, and the cart pulls back on the horse with an equal force. How can the horse and cart speed up?
\choices{The horse pulls harder than the cart pulls back.}{The ground pushes forward on the horse's hooves more than friction holds back the cart.}{The third law does not apply to animals.}{The cart pulls harder than the horse.}
\answer{B}{For the horse-cart system, the horse-cart forces are internal. The external forward force is the friction of the ground on the hooves.}
\end{exo}
\Needspace{12\baselineskip}
\begin{exo}[Free response]{}
Block A ($2\kg$) and block B ($3\kg$) are side by side on a frictionless floor. A horizontal force of $10\N$ pushes A to the right, and A pushes B.
\begin{enumerate}[label=(\alph*)]
\item Find the acceleration of the blocks.
\item Find the force of A on B.
\item Find the force of B on A (magnitude and direction). Which law did you use?
\item Check that the net force on A gives its acceleration.
\end{enumerate}
\work{4.5cm}
\begin{solution}
(a) $a=10/5=2\mss$. \;(b) $F_{A\text{ on }B}=m_Ba=3(2)=6\N$ to the right. \;(c) $6\N$ to the left (Newton's third law). \;(d) $10-6=4\N=2\kg\times2\mss$. \checkmark
\end{solution}
\end{exo}

%%DOC 2.4A | Newton's First Law
\begin{exo}{}
An object moves at constant velocity. What is the net force on it?
\choices{Zero}{In the direction of motion}{Opposite to the motion}{Increasing}
\answer{A}{Constant velocity means zero acceleration, so the net force is zero.}
\end{exo}
\begin{exo}{}
A puck slides on frictionless ice with no net force acting on it. It will
\choices{slow down}{stop}{keep moving at constant velocity}{speed up}
\answer{C}{First law: no net force, no change in velocity.}
\end{exo}
\begin{exo}{}
A $500\kg$ elevator moves upward at a constant speed. What is the tension in the cable?
\choices{$0$}{Less than $5000\N$}{More than $5000\N$}{$5000\N$}
\answer{D}{Equilibrium: $F_T=mg=5000\N$.}
\end{exo}
\begin{exo}{}
The inertia of an object depends on its
\choices{velocity}{mass}{position}{applied force}
\answer{B}{Mass measures inertia.}
\end{exo}
\begin{exo}{}
A passenger lurches forward when a bus brakes suddenly. Why?
\choices{A forward force pushes him.}{His body keeps moving while the bus slows down.}{Gravity pulls him forward.}{The seat pushes him forward.}
\answer{B}{By inertia, his body tends to keep its velocity; the bus slows under him.}
\end{exo}
\begin{exo}{}
A $100\N$ sign hangs from two ropes, each at $30^\circ$ above the horizontal. What is the tension in each rope?
\choices{$50\N$}{$57.7\N$}{$100\N$}{$200\N$}
\answer{C}{Vertical balance: $2T\sin30^\circ=100$, so $T=100\N$.}
\end{exo}
\begin{exo}{}
A box is pushed to the right with $30\N$ and moves at constant velocity. What is the friction force?
\choices{$30\N$ to the left}{$0$}{Less than $30\N$}{$30\N$ to the right}
\answer{A}{Constant velocity: the forces balance.}
\end{exo}
\begin{exo}{}
Which object is in equilibrium?
\choices{A ball at the top of its flight}{A car turning a corner at constant speed}{A car moving at constant speed on a straight, level road}{A falling apple}
\answer{C}{Constant speed in a straight line means constant velocity. The others accelerate.}
\end{exo}
\Needspace{12\baselineskip}
\begin{exo}[Free response]{}
A $10\kg$ box is held at rest on a frictionless $30^\circ$ incline by a rope parallel to the incline.
\begin{enumerate}[label=(\alph*)]
\item Draw a free-body diagram of the box.
\item Find the tension in the rope.
\item Find the normal force.
\item The rope now pulls with $60\N$. Is the box in equilibrium? Find its acceleration.
\end{enumerate}
\work{5cm}
\begin{solution}
(a) Weight down, normal force perpendicular to the incline, tension up the incline. \;(b) $F_T=mg\sin30^\circ=50\N$. \;(c) $F_N=mg\cos30^\circ\approx86.6\N$. \;(d) No: net force $60-50=10\N$ up the incline, $a=1\mss$ up the incline.
\end{solution}
\end{exo}

%%DOC 2.5A | Newton's Second Law
\begin{exo}{}
A net force of $20\N$ acts on a $4\kg$ object. What is its acceleration?
\choices{$80\mss$}{$5\mss$}{$0.2\mss$}{$24\mss$}
\answer{B}{$a=20/4=5\mss$.}
\end{exo}
\begin{exo}{}
A $1500\kg$ car accelerates at $2\mss$. What is the net force on it?
\choices{$750\N$}{$1502\N$}{$3000\N$}{$3\N$}
\answer{C}{$F=ma=1500(2)=3000\N$.}
\end{exo}
\begin{exo}{}
The same net force acts on an object of mass $m$ and on an object of mass $2m$. The acceleration of the second object is
\choices{twice as large}{half as large}{the same}{four times as large}
\answer{B}{$a=F/m$: doubling the mass halves the acceleration.}
\end{exo}
\begin{exo}{}
A $5\kg$ box is pushed with $40\N$ while friction is $15\N$. What is its acceleration?
\choices{$8\mss$}{$11\mss$}{$3\mss$}{$5\mss$}
\answer{D}{$a=(40-15)/5=5\mss$.}
\end{exo}
\begin{exo}{}
A $50\kg$ person stands on a scale in an elevator accelerating upward at $2\mss$. What does the scale read?
\choices{$400\N$}{$500\N$}{$600\N$}{$100\N$}
\answer{C}{$F_N=m(g+a)=50(12)=600\N$.}
\end{exo}
\begin{exo}{}
The same elevator now accelerates downward at $2\mss$. What does the scale read?
\choices{$400\N$}{$500\N$}{$600\N$}{$0$}
\answer{A}{$F_N=m(g-a)=50(8)=400\N$.}
\end{exo}
\begin{exo}{}
A $3\kg$ block on a frictionless table is tied by a string over a pulley to a hanging $1\kg$ block. What is the acceleration?
\choices{$10\mss$}{$2.5\mss$}{$3.3\mss$}{$7.5\mss$}
\answer{B}{System: $a=\dfrac{m_2g}{m_1+m_2}=\dfrac{10}{4}=2.5\mss$.}
\end{exo}
\begin{exo}{}
What is the tension in the string in the previous question?
\choices{$10\N$}{$2.5\N$}{$7.5\N$}{$30\N$}
\answer{C}{On the $3\kg$ block: $F_T=3(2.5)=7.5\N$ (less than the $10\N$ weight of the hanging block, which accelerates down).}
\end{exo}
\Needspace{12\baselineskip}
\begin{exo}[Free response]{}
A $3\kg$ block and a $2\kg$ block hang on either side of a light string over a frictionless pulley (an Atwood machine). They are released from rest.
\begin{enumerate}[label=(\alph*)]
\item Draw a free-body diagram for each block.
\item Find the acceleration of the blocks.
\item Find the tension in the string.
\item Both masses are doubled. How do the acceleration and the tension change? Explain.
\end{enumerate}
\work{5cm}
\begin{solution}
(a) Each block: tension up, weight down; the $3\kg$ block's weight arrow is longer than the tension, the $2\kg$ block's is shorter. \;(b) $a=\frac{(3-2)(10)}{5}=2\mss$. \;(c) On the $2\kg$ block: $F_T-20=2(2)$, $F_T=24\N$. \;(d) $a=\frac{(m_1-m_2)g}{m_1+m_2}$ is unchanged; the tension doubles to $48\N$.
\end{solution}
\end{exo}

%%DOC 2.6A | Gravitational Force
\begin{exo}{}
The distance between two objects is doubled. The gravitational force between them becomes
\choices{half as large}{twice as large}{one quarter as large}{the same}
\answer{C}{Inverse square: $1/2^2=1/4$.}
\end{exo}
\begin{exo}{}
The mass of one of the objects is tripled. The gravitational force becomes
\choices{three times as large}{nine times as large}{one third as large}{the same}
\answer{A}{$F_g$ is proportional to each mass.}
\end{exo}
\begin{exo}{}
A $70\kg$ astronaut travels to the Moon ($g=1.6\,\mathrm{N/kg}$). What is her mass on the Moon?
\choices{$70\kg$}{$11.2\kg$}{$112\kg$}{$0$}
\answer{A}{Mass does not depend on location.}
\end{exo}
\begin{exo}{}
What is the weight of the same astronaut on the Moon?
\choices{$700\N$}{$70\N$}{$112\N$}{$43.8\N$}
\answer{C}{$F_g=mg=70(1.6)=112\N$.}
\end{exo}
\begin{exo}{}
What is the gravitational field at a height above Earth's surface equal to Earth's radius?
\choices{$5\,\mathrm{N/kg}$}{$2.5\,\mathrm{N/kg}$}{$10\,\mathrm{N/kg}$}{$1.25\,\mathrm{N/kg}$}
\answer{B}{The distance to Earth's center doubles: $10/4=2.5\,\mathrm{N/kg}$.}
\end{exo}
\begin{exo}{}
Astronauts in an orbiting space station feel weightless because
\choices{there is no gravity in space}{they are in free fall with the station}{they are far from Earth}{the station blocks gravity}
\answer{B}{Gravity still acts on them, but everything falls together, so nothing pushes on them.}
\end{exo}
\begin{exo}{}
What is the gravitational field strength near Earth's surface?
\choices{$10\,\mathrm{kg/N}$}{$6.67\times10^{-11}\,\mathrm{N/kg}$}{$1\,\mathrm{N/kg}$}{$10\,\mathrm{N/kg}$}
\answer{D}{$g\approx10\,\mathrm{N/kg}$, the same as $10\mss$.}
\end{exo}
\begin{exo}{}
What is the gravitational force between two $1000\kg$ masses whose centers are $1\m$ apart? ($G=6.67\times10^{-11}\,\mathrm{N\,m^2/kg^2}$)
\choices{$6.67\times10^{-5}\N$}{$6.67\times10^{-11}\N$}{$6.67\times10^{-8}\N$}{$6.67\N$}
\answer{A}{$F=G\frac{(10^3)(10^3)}{1^2}=6.67\times10^{-5}\N$: tiny, which is why we do not notice it.}
\end{exo}
\Needspace{12\baselineskip}
\begin{exo}[Free response]{}
Planet X has twice the mass of Earth and twice its radius. On Earth, $g=10\,\mathrm{N/kg}$.
\begin{enumerate}[label=(\alph*)]
\item Find the gravitational field at the surface of planet X. Show your reasoning with $g=GM/R^2$.
\item Find the weight of a $40\kg$ student on planet X.
\item What is the student's mass on planet X?
\end{enumerate}
\work{4.5cm}
\begin{solution}
(a) $g_X=\frac{G(2M)}{(2R)^2}=\frac{2}{4}\,g=5\,\mathrm{N/kg}$. \;(b) $40(5)=200\N$. \;(c) $40\kg$: mass does not change.
\end{solution}
\end{exo}

%%DOC 2.7A | Kinetic and Static Friction
\begin{exo}{}
A $10\kg$ box slides on a level floor with $\mu_k=0.3$. What is the kinetic friction force?
\choices{$3\N$}{$30\N$}{$300\N$}{$0.3\N$}
\answer{B}{$f_k=\mu_kmg=0.3(100)=30\N$.}
\end{exo}
\begin{exo}{}
The same box is at rest, with $\mu_s=0.5$. What is the largest push it can take without starting to slide?
\choices{$50\N$}{$30\N$}{$100\N$}{$5\N$}
\answer{A}{$f_{s,\max}=\mu_smg=0.5(100)=50\N$.}
\end{exo}
\begin{exo}{}
You push the box at rest with $20\N$. What is the static friction force?
\choices{$50\N$}{$0$}{$20\N$}{$30\N$}
\answer{C}{Static friction matches the push (up to $50\N$): $20\N$, and the box stays at rest.}
\end{exo}
\begin{exo}{}
The box slides and you push it with $40\N$ ($f_k=30\N$). What is its acceleration?
\choices{$4\mss$}{$1\mss$}{$7\mss$}{$3\mss$}
\answer{B}{$a=(40-30)/10=1\mss$.}
\end{exo}
\begin{exo}{}
A box slides on a level floor with no push, and $\mu_k=0.2$. What is the magnitude of its acceleration?
\choices{$0.2\mss$}{$20\mss$}{$2\mss$}{$10\mss$}
\answer{C}{$a=f_k/m=\mu_kg=0.2(10)=2\mss$ (it slows down).}
\end{exo}
\begin{exo}{}
A brick slides on its large face, then on its small face. In the model used in AP Physics 1, the kinetic friction is
\choices{larger on the large face}{smaller on the large face}{zero on the small face}{the same}
\answer{D}{$f_k=\mu_kF_N$ does not depend on the contact area.}
\end{exo}
\begin{exo}{}
A $2\kg$ block rests on a $30^\circ$ incline. What is the static friction force on it?
\choices{$10\N$ up the incline}{$17.3\N$ up the incline}{$10\N$ down the incline}{$20\N$ up the incline}
\answer{A}{It balances the component of gravity along the incline: $mg\sin30^\circ=20(0.5)=10\N$.}
\end{exo}
\begin{exo}{}
What is the normal force on the block in the previous question?
\choices{$20\N$}{$10\N$}{$17.3\N$}{$0$}
\answer{C}{$F_N=mg\cos30^\circ=20(0.866)\approx17.3\N$: less than the weight.}
\end{exo}
\Needspace{12\baselineskip}
\begin{exo}[Free response]{}
A $4\kg$ box is pulled across a level floor by a horizontal rope with a tension of $25\N$. The coefficient of kinetic friction is $0.4$.
\begin{enumerate}[label=(\alph*)]
\item Draw a free-body diagram of the box.
\item Find the friction force.
\item Find the acceleration of the box.
\item The rope is now pulled at an angle above the horizontal, with the same tension. Does the friction force increase, decrease or stay the same? Explain without calculating.
\end{enumerate}
\work{5cm}
\begin{solution}
(a) Weight down, normal up, tension right, friction left. \;(b) $f_k=0.4(40)=16\N$. \;(c) $a=(25-16)/4=2.25\mss$. \;(d) It decreases: the rope now also pulls up, so the floor pushes up less ($F_N<mg$), and $f_k=\mu_kF_N$ is smaller.
\end{solution}
\end{exo}

%%DOC 2.8A | Spring Forces
\begin{exo}{}
A spring with $k=200\Nm$ is stretched by $0.05\m$. What force does it exert?
\choices{$4000\N$}{$10\N$}{$40\N$}{$0.00025\N$}
\answer{B}{$F=k\Delta x=200(0.05)=10\N$.}
\end{exo}
\begin{exo}{}
A force of $30\N$ stretches a spring by $0.15\m$. What is its spring constant?
\choices{$4.5\Nm$}{$0.005\Nm$}{$200\Nm$}{$20\Nm$}
\answer{C}{$k=F/\Delta x=30/0.15=200\Nm$.}
\end{exo}
\begin{exo}{}
In which direction does a spring force act on the object attached to it?
\choices{In the direction of the stretch}{Back toward the spring's natural length}{Always upward}{Always downward}
\answer{B}{It is a restoring force: $F=-k\Delta x$.}
\end{exo}
\begin{exo}{}
A $2\kg$ mass hangs at rest from a spring with $k=400\Nm$. How much is the spring stretched?
\choices{$0.05\m$}{$0.5\m$}{$5\m$}{$0.005\m$}
\answer{A}{$k\Delta x=mg\Rightarrow\Delta x=20/400=0.05\m$.}
\end{exo}
\begin{exo}{}
On a graph of spring force versus stretch, the slope represents
\choices{the spring constant $k$}{$1/k$}{the mass}{the energy}
\answer{A}{$F=k\Delta x$: a straight line of slope $k$.}
\end{exo}
\begin{exo}{}
The stretch of a spring is doubled. The spring force
\choices{stays the same}{is halved}{doubles}{quadruples}
\answer{C}{The force is proportional to the stretch.}
\end{exo}
\begin{exo}{}
A spring is compressed by $0.1\m$ with a force of $15\N$. What force does it exert when it is stretched by $0.2\m$?
\choices{$15\N$}{$7.5\N$}{$60\N$}{$30\N$}
\answer{D}{$k=150\Nm$; $150(0.2)=30\N$.}
\end{exo}
\begin{exo}{}
A stiffer spring has
\choices{a smaller spring constant}{a larger spring constant}{the same spring constant}{no spring constant}
\answer{B}{A stiffer spring needs more force for the same stretch: larger $k$.}
\end{exo}
\Needspace{16\baselineskip}
\begin{exo}[Free response \textbullet\ Experimental design and analysis]{}
Students hang masses from a spring and measure the stretch.
\begin{center}\begin{tabular}{c|cccc}
$m$ (kg) & $0.10$ & $0.20$ & $0.30$ & $0.40$\\\hline
stretch (m) & $0.020$ & $0.041$ & $0.059$ & $0.080$
\end{tabular}\end{center}
\begin{enumerate}[label=(\alph*)]
\item Find the spring force for each mass.
\item What should they plot to find $k$ from a slope? Find $k$.
\item Predict the stretch for a $0.25\kg$ mass.
\end{enumerate}
\work{5cm}
\begin{solution}
(a) $F=mg$: $1$, $2$, $3$, $4\N$. \;(b) Force (vertical axis) vs.\ stretch: slope $=k\approx\frac{4-1}{0.080-0.020}=50\Nm$. \;(c) $\Delta x=2.5/50=0.05\m$.
\end{solution}
\end{exo}

%%DOC 2.9A | Circular Motion
\begin{exo}{}
In uniform circular motion, which quantity stays constant?
\choices{The velocity}{The speed}{The acceleration vector}{The direction of motion}
\answer{B}{Only the magnitude of the velocity is constant; its direction keeps changing.}
\end{exo}
\begin{exo}{}
In uniform circular motion, the acceleration points
\choices{along the tangent}{away from the center}{toward the center}{nowhere: it is zero}
\answer{C}{Centripetal means ``toward the center''.}
\end{exo}
\begin{exo}{}
A car moves at $10\ms$ around a curve of radius $20\m$. What is its centripetal acceleration?
\choices{$0.5\mss$}{$2\mss$}{$200\mss$}{$5\mss$}
\answer{D}{$a_c=v^2/r=100/20=5\mss$.}
\end{exo}
\begin{exo}{}
A $0.2\kg$ ball on a string moves in a horizontal circle of radius $0.5\m$ at $2\ms$ (ignore gravity). What is the tension?
\choices{$0.8\N$}{$1.6\N$}{$0.4\N$}{$8\N$}
\answer{B}{$a_c=4/0.5=8\mss$; $F_T=0.2(8)=1.6\N$.}
\end{exo}
\begin{exo}{}
The string breaks. The ball then moves
\choices{toward the center}{straight outward along the radius}{in a straight line along the tangent}{in the same circle}
\answer{C}{With no net force, it keeps the velocity it had: tangent to the circle.}
\end{exo}
\begin{exo}{}
A car goes around a flat curve at constant speed. Which force provides the centripetal force?
\choices{Static friction from the road}{The normal force}{Gravity}{The engine}
\answer{A}{The tires grip the road sideways (no sliding): static friction points toward the center.}
\end{exo}
\begin{exo}{}
The speed of an object in a circle of fixed radius doubles. The net force needed becomes
\choices{twice as large}{four times as large}{half as large}{the same}
\answer{B}{$F=mv^2/r$: $2^2=4$.}
\end{exo}
\begin{exo}{}
A ball goes around a circle of radius $2\m$ once every $2\s$. What is its speed?
\choices{About $6.3\ms$}{About $3.1\ms$}{About $12.6\ms$}{$1\ms$}
\answer{A}{$v=2\pi r/T=2\pi(2)/2=2\pi\approx6.3\ms$.}
\end{exo}
\Needspace{12\baselineskip}
\begin{exo}[Free response]{}
A $1000\kg$ car drives over the top of a hill shaped like a circular arc of radius $40\m$, at $10\ms$.
\begin{enumerate}[label=(\alph*)]
\item Draw a free-body diagram of the car at the top.
\item Find the centripetal acceleration and the normal force on the car.
\item At what speed would the car lose contact with the road at the top? Explain.
\end{enumerate}
\work{5cm}
\begin{solution}
(a) Weight down (longer arrow), normal force up. \;(b) $a_c=100/40=2.5\mss$ downward; $mg-F_N=ma_c\Rightarrow F_N=1000(10-2.5)=7500\N$. \;(c) When $F_N=0$: $mg=mv^2/r$, $v=\sqrt{gr}=\sqrt{400}=20\ms$.
\end{solution}
\end{exo}

%%DOC SC2.1 | Self Check 2.1 \textbullet\ Topics 2.1 to 2.5
\begin{exo}{}
A $2\kg$ mass is at $x=1\m$ and another $2\kg$ mass is at $x=5\m$. Where is their center of mass?
\choices{$x=3\m$}{$x=6\m$}{$x=2\m$}{$x=4\m$}
\answer{A}{Equal masses: halfway, $3\m$.}
\end{exo}
\begin{exo}{}
Earth pulls a ball downward. What is the third-law partner of this force?
\choices{The normal force on the ball}{The ball pulling Earth upward}{Air resistance}{The ball pushing on the ground}
\answer{B}{Same pair of objects, same type of force, opposite direction.}
\end{exo}
\begin{exo}{}
The net force on a moving object is zero. The object
\choices{slows down}{stops}{keeps a constant velocity}{speeds up}
\answer{C}{First law.}
\end{exo}
\begin{exo}{}
A net force of $18\N$ acts on a $6\kg$ object. What is its acceleration?
\choices{$108\mss$}{$0.33\mss$}{$12\mss$}{$3\mss$}
\answer{D}{$18/6=3\mss$.}
\end{exo}
\begin{exo}{}
A $70\kg$ person rides an elevator moving down at a constant speed. What does a scale under her read?
\choices{$0$}{Less than $700\N$}{$700\N$}{More than $700\N$}
\answer{C}{Constant velocity: no acceleration, so $F_N=mg$.}
\end{exo}
\begin{exo}{}
For a book at rest on a table, why is the normal force equal to the weight?
\choices{They are a third-law pair.}{The net force on the book is zero (first law).}{The normal force always equals the weight.}{Gravity is weak.}
\answer{B}{They act on the same object and the book is in equilibrium. (They are not a third-law pair, and $F_N\neq mg$ in an accelerating elevator or on an incline.)}
\end{exo}
\Needspace{12\baselineskip}
\begin{exo}[Free response]{}
A $2\kg$ block on a frictionless table is tied by a string over a pulley to a hanging $3\kg$ block.
\begin{enumerate}[label=(\alph*)]
\item Draw a free-body diagram for each block.
\item Find the acceleration.
\item Find the tension in the string.
\end{enumerate}
\work{4.5cm}
\begin{solution}
(a) Table block: weight, normal, tension toward the pulley. Hanging block: tension up, weight down (longer). \;(b) $a=\frac{30}{5}=6\mss$. \;(c) $F_T=2(6)=12\N$ (check: $30-12=18=3\times6$).
\end{solution}
\end{exo}

%%DOC SC2.2 | Self Check 2.2 \textbullet\ Topics 2.6 to 2.9
\begin{exo}{}
Two objects $3\m$ apart attract each other with a force of $36\N$. What is the force when they are $6\m$ apart?
\choices{$18\N$}{$9\N$}{$72\N$}{$4\N$}
\answer{B}{Distance doubled: force divided by $4$.}
\end{exo}
\begin{exo}{}
A planet has the same radius as Earth but twice its mass. What is the gravitational field at its surface?
\choices{$5\,\mathrm{N/kg}$}{$10\,\mathrm{N/kg}$}{$20\,\mathrm{N/kg}$}{$40\,\mathrm{N/kg}$}
\answer{C}{$g=GM/R^2$ doubles: $20\,\mathrm{N/kg}$.}
\end{exo}
\begin{exo}{}
A $5\kg$ box slides on a level floor with $\mu_k=0.4$. What is the friction force?
\choices{$20\N$}{$2\N$}{$50\N$}{$12.5\N$}
\answer{A}{$0.4(50)=20\N$.}
\end{exo}
\begin{exo}{}
A box at rest can take up to $30\N$ of static friction. You push it with $10\N$. What is the friction force?
\choices{$30\N$}{$0$}{$20\N$}{$10\N$}
\answer{D}{Static friction is only as large as needed: $10\N$.}
\end{exo}
\begin{exo}{}
A spring with $k=50\Nm$ is compressed by $0.2\m$. What force does it exert?
\choices{$250\N$}{$10\N$}{$2.5\N$}{$0.004\N$}
\answer{B}{$50(0.2)=10\N$.}
\end{exo}
\begin{exo}{}
A ball on a string moves in a vertical circle. At the lowest point, how does the tension compare with the ball's weight?
\choices{$F_T=mg$}{$F_T<mg$}{$F_T>mg$}{$F_T=0$}
\answer{C}{The net force must point up toward the center: $F_T-mg=mv^2/r>0$.}
\end{exo}
\Needspace{12\baselineskip}
\begin{exo}[Free response]{}
A car drives around a flat curve of radius $50\m$. The coefficient of static friction between the tires and the road is $0.8$.
\begin{enumerate}[label=(\alph*)]
\item Draw a free-body diagram of the car (seen from behind).
\item Find the maximum speed for which the car does not skid.
\item When the road is wet, $\mu_s$ is halved. Find the new maximum speed.
\end{enumerate}
\work{4.5cm}
\begin{solution}
(a) Weight down, normal up, static friction horizontal toward the center of the curve. \;(b) $\mu_smg=mv^2/r\Rightarrow v=\sqrt{\mu_sgr}=\sqrt{0.8(10)(50)}=20\ms$. \;(c) $v=\sqrt{0.4(10)(50)}\approx14.1\ms$ ($20/\sqrt2$). The mass does not matter.
\end{solution}
\end{exo}

%%DOC PRETEST2 | Unit 2 Test \textbullet\ Force and Translational Dynamics
\partheading{Part A \textbullet\ Multiple choice}
\begin{exo}{}
A $1\kg$ mass is at $x=0$ and a $4\kg$ mass is at $x=5\m$. Where is their center of mass?
\choices{$x=2.5\m$}{$x=4\m$}{$x=1\m$}{$x=5\m$}
\answer{B}{$\frac{0+20}{5}=4\m$.}
\end{exo}
\begin{exo}{}
Which of these is a contact force?
\choices{Friction}{Gravity}{The magnetic force}{The electric force}
\answer{A}{Friction needs two surfaces in contact.}
\end{exo}
\begin{exo}{}
A bug hits the windshield of a moving car. Compare the force of the bug on the car with the force of the car on the bug.
\choices{The bug's force is larger.}{The car's force is larger.}{They are equal in magnitude.}{The bug's force is zero.}
\answer{C}{Third-law pair. The bug's acceleration is much larger because its mass is tiny.}
\end{exo}
\begin{exo}{}
Three forces act on an object moving at constant velocity. The net force is
\choices{equal to the largest force}{in the direction of motion}{impossible to know}{zero}
\answer{D}{Constant velocity: equilibrium.}
\end{exo}
\begin{exo}{}
A $2\kg$ object is pushed by $10\N$ to the right and $4\N$ to the left. What is its acceleration?
\choices{$7\mss$ right}{$3\mss$ right}{$3\mss$ left}{$5\mss$ right}
\answer{B}{Net $6\N$ right; $a=6/2=3\mss$ right.}
\end{exo}
\begin{exo}{}
An $80\kg$ person stands on a scale in an elevator accelerating upward at $1\mss$. What does the scale read?
\choices{$800\N$}{$720\N$}{$880\N$}{$80\N$}
\answer{C}{$80(10+1)=880\N$.}
\end{exo}
\begin{exo}{}
What is the weight of a $50\kg$ mass on a planet where $g=4\,\mathrm{N/kg}$?
\choices{$500\N$}{$200\N$}{$12.5\N$}{$50\N$}
\answer{B}{$50(4)=200\N$.}
\end{exo}
\begin{exo}{}
The distance between two objects is tripled. The gravitational force between them becomes
\choices{one third as large}{three times as large}{nine times as large}{one ninth as large}
\answer{D}{Inverse square: $1/3^2=1/9$.}
\end{exo}
\begin{exo}{}
An $8\kg$ box slides on a level floor ($\mu_k=0.25$) while pushed with $30\N$. What is its acceleration?
\choices{$1.25\mss$}{$3.75\mss$}{$2.5\mss$}{$0$}
\answer{A}{$f_k=0.25(80)=20\N$; $a=(30-20)/8=1.25\mss$.}
\end{exo}
\begin{exo}{}
The maximum static friction on a box at rest is $12\N$. You push it with $8\N$. What is the friction force?
\choices{$12\N$}{$4\N$}{$8\N$}{$0$}
\answer{C}{Static friction equals the push: $8\N$. The box does not move.}
\end{exo}
\begin{exo}{}
A $3\kg$ mass hangs at rest from a spring with $k=300\Nm$. How much is the spring stretched?
\choices{$1\m$}{$0.1\m$}{$0.01\m$}{$10\m$}
\answer{B}{$\Delta x=mg/k=30/300=0.1\m$.}
\end{exo}
\begin{exo}{}
In uniform circular motion, the net force on the object points
\choices{toward the center}{along the velocity}{away from the center}{nowhere: it is zero}
\answer{A}{The net force is centripetal.}
\end{exo}
\begin{exo}{}
A $1200\kg$ car rounds a curve of radius $45\m$ at $15\ms$. What net force is needed?
\choices{$400\N$}{$18\,000\N$}{$600\N$}{$6000\N$}
\answer{D}{$mv^2/r=1200(225)/45=6000\N$.}
\end{exo}
\begin{exo}{}
A block slides down a frictionless $30^\circ$ incline. What is its acceleration?
\choices{$10\mss$}{$8.7\mss$}{$5\mss$}{$0$}
\answer{C}{Along the incline: $a=g\sin30^\circ=5\mss$.}
\end{exo}
\Needspace{24\baselineskip}
\partheading{Part B \textbullet\ Free response}
\begin{exo}[Free response \textbullet\ Mathematical routines]{}
A $4\kg$ block on a rough table ($\mu_k=0.25$) is tied by a light string over a light, frictionless pulley to a hanging $2\kg$ block. The system is released from rest.
\begin{enumerate}[label=(\alph*)]
\item Draw a free-body diagram for each block.
\item Find the acceleration of the blocks.
\item Find the tension in the string.
\item If the table were frictionless, would the tension be larger, smaller or the same? Justify.
\end{enumerate}
\work{6cm}
\begin{solution}
(a) Table block: weight, normal, tension toward the pulley, friction away from it. Hanging block: tension up, weight down. \;(b) System: $a=\frac{20-0.25(40)}{6}=\frac{10}{6}\approx1.67\mss$. \;(c) Hanging block: $20-F_T=2(1.67)$, $F_T\approx16.7\N$. \;(d) Smaller: without friction $a=20/6\approx3.33\mss$, and $F_T=2(10-3.33)\approx13.3\N$. The hanging block accelerates more, so the string holds it back less.
\end{solution}
\end{exo}
\Needspace{16\baselineskip}
\begin{exo}[Free response \textbullet\ Qualitative and quantitative translation]{}
A $0.5\kg$ ball on a string of length $0.8\m$ moves in a vertical circle.
\begin{enumerate}[label=(\alph*)]
\item Without calculating, explain whether the tension is larger at the top or at the bottom of the circle.
\item At the top, the speed is $4\ms$. Find the tension.
\item At the bottom, the speed is $6\ms$. Find the tension.
\item Is the result of (b) and (c) consistent with your answer to (a)?
\end{enumerate}
\work{6cm}
\begin{solution}
(a) At the bottom: there, tension must overcome gravity to point the net force up toward the center, and the ball is also faster. At the top, gravity already points toward the center and helps. \;(b) $F_T+mg=mv^2/r$: $F_T=0.5(16/0.8)-5=10-5=5\N$. \;(c) $F_T-mg=mv^2/r$: $F_T=5+0.5(36/0.8)=5+22.5=27.5\N$. \;(d) Yes: $27.5\N>5\N$.
\end{solution}
\end{exo}
